\section{Zadanie 2}
Czy istnieje nieskończona przeliczalna \(\sigma\)-algebra?

\subsection*{Rozwiązanie}
Nie istnieje nieskończona przeliczalna \(\sigma\)-algebra. Pokazmy sobie to na atomach.
Atom w \(\sigma\)-algebrze to taki zbiór \(A\) należący do tej \(\sigma\)-algebry, że dla dowolnego podzbioru \(B \subseteq A\) zachodzi \(B = \emptyset\) lub \(B = A\). Oznacza to, że atom nie może być podzielony na mniejsze zbiory należące do tej samej \(\sigma\)-algebry.

Widzimy ze jesli mamy skonczoną liczbę atomów to \(\sigma\)-algebra jest skończona, bo każdy zbiór w \(\sigma\)-algebrze można przedstawić jako sumę atomów. Czyli mamy \(2^n\) podzbiorów dla \(n\) atomów.

To teraz załózmy ze mamy nieskończenie i przeliczalnie wiele atomów. I zastanówmy się ile mamy podzbiorów.
No mamy ich \(|\mathbb{P}(A)| = 2^A = 2^{\aleph_0}\) poniewaz rozwazamy \(A\) które jest nieskończone i przeliczalne.
A z MFI wiemy ze \(2^{\aleph_0}\) jest nieprzeliczalne bo jest równe na moc continuum.

\subsection*{Odpowiedź}
\[
	\boxed{\text{Nie}}
\]
